
%----------------- LISTA DE ABREVIATURAS E SIGLAS--------------------------------------
\siglalista{ABNT}{Associação Brasileira de Normas Técnicas}

%Para usar uma dada sigla ABC ao longo do texto, use \glsxtrfull{ABC} se quiser apresentar a sigla e sua definição.
%Se quiser apresentar apenas a sigla, use \gls{ABC}.






%-----------------SÍMBOLOS---------------------------------------------------------------
\simbololista{G}{\ensuremath{G}}{Resultante do peso próprio (ação permanente)}
\simbololista{gamma_med}{\ensuremath{\gamma_{med}}}{Valor médio do peso específico do material}
\simbololista{V}{\ensuremath{V}}{Volume do elemento estrutural}
\simbololista{F_d}{\ensuremath{F_d}}{Valor de cálculo da ação combinada}
\simbololista{gamma_gi}{\ensuremath{\gamma_{gi}}}{Coeficientes de majoração das ações permanentes}
\simbololista{F_Gi,k}{\ensuremath{F_{Gi,k}}}{Valores característicos das ações permanentes}
\simbololista{gamma_q1}{\ensuremath{\gamma_{q1}}}{Coeficiente de majoração da ação variável principal}
\simbololista{gamma_qj}{\ensuremath{\gamma_{qj}}}{Coeficientes de majoração das ações variáveis secundárias}
\simbololista{F_Q1,k}{\ensuremath{F_{Q1,k}}}{Valor característico da ação variável principal}
\simbololista{psi_0j}{\ensuremath{\psi_{0j}}}{Fatores de combinação das ações variáveis secundárias}
\simbololista{F_Qj,k}{\ensuremath{F_{Qj,k}}}{Valores característicos das ações variáveis secundárias}
\simbololista{m}{\ensuremath{m}}{Número de ações permanentes}
\simbololista{n}{\ensuremath{n}}{Número de ações variáveis secundárias}
\simbololista{gamma_qw}{\ensuremath{\gamma_{qw}}}{Coeficiente de majoração da ação do vento}
\simbololista{gamma_q2}{\ensuremath{\gamma_{q2}}}{Coeficiente de majoração da ação variável secundária}
\simbololista{psi_0,q2}{\ensuremath{\psi_{0,q2}}}{Fator de combinação da ação variável secundária}
\simbololista{X_d}{\ensuremath{X_d}}{Valor de cálculo da resistência da madeira}
\simbololista{X_k}{\ensuremath{X_k}}{Valor característico da resistência da madeira}
\simbololista{k_mod}{\ensuremath{k_{mod}}}{Coeficiente de modificação da madeira}
\simbololista{gamma_w}{\ensuremath{\gamma_w}}{Coeficiente de ponderação da resistência}
\simbololista{k_mod,1}{\ensuremath{k_{mod,1}}}{Fator de modificação associado à classe de carregamento}
\simbololista{k_mod,2}{\ensuremath{k_{mod,2}}}{Fator de modificação associado à classe de umidade}
\simbololista{E_0,ef}{\ensuremath{E_{0,ef}}}{Valor efetivo do módulo de elasticidade}
\simbololista{E_0,med}{\ensuremath{E_{0,med}}}{Valor médio do módulo de elasticidade}
\simbololista{sigma_c0,d}{\ensuremath{\sigma_{c0,d}}}{Tensão normal de cálculo na compressão paralela às fibras}
\simbololista{M_sd}{\ensuremath{M_{sd}}}{Momento fletor solicitante de cálculo}
\simbololista{W_c}{\ensuremath{W_c}}{Módulo de resistência da seção transversal comprimida}
\simbololista{f_b,d}{\ensuremath{f_{b,d}}}{Resistência de cálculo à flexão}
\simbololista{sigma_r0,d}{\ensuremath{\sigma_{r0,d}}}{Tensão normal de cálculo na tração paralela às fibras}
\simbololista{W_t}{\ensuremath{W_t}}{Módulo de resistência da seção transversal tracionada}
\simbololista{I}{\ensuremath{I}}{Momento de inércia da seção transversal}
\simbololista{y_c}{\ensuremath{y_c}}{Distância da linha neutra ao bordo comprimido}
\simbololista{y_t}{\ensuremath{y_t}}{Distância da linha neutra ao bordo tracionado}
\simbololista{b}{\ensuremath{b}}{Largura da seção transversal retangular}
\simbololista{h}{\ensuremath{h}}{Altura da seção transversal retangular}
\simbololista{R}{\ensuremath{R}}{Raio da seção transversal circular}
\simbololista{D}{\ensuremath{D}}{Diâmetro da seção transversal circular}
\simbololista{tau_d}{\ensuremath{\tau_d}}{Valor de cálculo da tensão de cisalhamento}
\simbololista{V_d}{\ensuremath{V_d}}{Esforço cortante de cálculo}
\simbololista{S}{\ensuremath{S}}{Momento estático da parte da seção transversal}
\simbololista{f_v0,d}{\ensuremath{f_{v0,d}}}{Resistência de cálculo ao cisalhamento paralelo às fibras}
\simbololista{A}{\ensuremath{A}}{Área da seção transversal}
\simbololista{tau_max}{\ensuremath{\tau_{max}}}{Tensão máxima de cisalhamento}
\simbololista{delta_inst}{\ensuremath{\delta_{inst}}}{Flecha instantânea}
\simbololista{delta_inst,G_i,k}{\ensuremath{\delta_{inst,G_i,k}}}{Flechas instantâneas das ações permanentes}
\simbololista{delta_inst,Q_1,k}{\ensuremath{\delta_{inst,Q_1,k}}}{Flecha da ação variável principal}
\simbololista{delta_inst,Q_j,k}{\ensuremath{\delta_{inst,Q_j,k}}}{Flechas das ações variáveis secundárias}
\simbololista{psi_1,j}{\ensuremath{\psi_{1,j}}}{Fator de redução para condição frequente}
\simbololista{delta_fin}{\ensuremath{\delta_{fin}}}{Flecha final}
\simbololista{delta_fin,G_i,k}{\ensuremath{\delta_{fin,G_i,k}}}{Flechas finais das ações permanentes}
\simbololista{delta_fin,Q_j,k}{\ensuremath{\delta_{fin,Q_j,k}}}{Flechas finais das ações variáveis}
\simbololista{delta_creep}{\ensuremath{\delta_{creep}}}{Flecha associada à fluência}
\simbololista{phi}{\ensuremath{\varphi}}{Coeficiente de fluência da madeira}
\simbololista{sigma_d,f}{\ensuremath{\sigma_{d,f}}}{Tensão máxima de cálculo de flexão}
\simbololista{f_d,f}{\ensuremath{f_{d,f}}}{Resistência de cálculo à flexão}
\simbololista{M}{\ensuremath{M}}{Momento fletor}
\simbololista{W}{\ensuremath{W}}{Módulo de resistência}
\simbololista{sigma_d,f,x}{\ensuremath{\sigma_{d,f,x}}}{Tensão de flexão no eixo x}
\simbololista{sigma_d,f,y}{\ensuremath{\sigma_{d,f,y}}}{Tensão de flexão no eixo y}
\simbololista{sigma_d,n}{\ensuremath{\sigma_{d,n}}}{Tensão normal de tração}
\simbololista{f_d,t}{\ensuremath{f_{d,t}}}{Resistência à tração}
\simbololista{f_d,c}{\ensuremath{f_{d,c}}}{Resistência à compressão}
\simbololista{Q_g,k}{\ensuremath{Q_{g,k}}}{Força cortante máxima das ações permanentes}
\simbololista{g}{\ensuremath{g}}{Carga permanente distribuída}
\simbololista{L}{\ensuremath{L}}{Vão estrutural}
\simbololista{Q_q,k}{\ensuremath{Q_{q,k}}}{Força cortante máxima reduzida}
\simbololista{P}{\ensuremath{P}}{Carga concentrada}
\simbololista{q}{\ensuremath{q}}{Carga distribuída}
\simbololista{a}{\ensuremath{a}}{Distância da carga ao apoio}
\simbololista{e}{\ensuremath{e}}{Excentricidade}
\simbololista{delta_g,k}{\ensuremath{\delta_{g,k}}}{Flecha máxima das ações permanentes}
\simbololista{E_M,ef}{\ensuremath{E_{M,ef}}}{Módulo efetivo de elasticidade}
\simbololista{delta_q,k}{\ensuremath{\delta_{q,k}}}{Flecha máxima da ação variável}
\simbololista{b_dist}{\ensuremath{b}}{Distância da carga ao apoio mais próximo}
\simbololista{M_r,k}{\ensuremath{M_{r,k}}}{Momento fletor característico máximo}
\simbololista{q_r}{\ensuremath{q_r}}{Carga distribuída equivalente}
\simbololista{a_r}{\ensuremath{a_r}}{Comprimento de distribuição}
\simbololista{L_r}{\ensuremath{L_r}}{Vão do tabuleiro}
\simbololista{k_def}{\ensuremath{k_{def}}}{Fator de fluência}
\simbololista{E_mean,fin}{\ensuremath{E_{mean,fin}}}{Módulo médio final de elasticidade}
\simbololista{G_mean,fin}{\ensuremath{G_{mean,fin}}}{Módulo médio final de cisalhamento}
\simbololista{K_ser,fin}{\ensuremath{K_{ser,fin}}}{Módulo final de deslizamento}
\simbololista{E_mean}{\ensuremath{E_{mean}}}{Módulo médio de elasticidade}
\simbololista{G_mean}{\ensuremath{G_{mean}}}{Módulo médio de cisalhamento}
\simbololista{K_ser}{\ensuremath{K_{ser}}}{Módulo de deslizamento}
\simbololista{psi_2}{\ensuremath{\psi_2}}{Fator quase-permanente}
\simbololista{k_def,1}{\ensuremath{k_{def,1}}}{Fator de fluência elemento 1}
\simbololista{k_def,2}{\ensuremath{k_{def,2}}}{Fator de fluência elemento 2}
\simbololista{D_min}{\ensuremath{D_{min}}}{Diâmetro mínimo da longarina}
\simbololista{M_d}{\ensuremath{M_d}}{Momento fletor de cálculo}
\simbololista{Q_d}{\ensuremath{Q_d}}{Esforço cortante de cálculo}
\simbololista{p_f}{\ensuremath{p_f}}{Probabilidade de falha}
\simbololista{X}{\ensuremath{\mathbf{X}}}{Vetor de variáveis aleatórias}
\simbololista{f_X}{\ensuremath{f_X}}{Função densidade de probabilidade}
\simbololista{G_lim}{\ensuremath{G}}{Equação do estado limite}
\simbololista{p_hat_f}{\ensuremath{\hat{p}_f}}{Estimativa da probabilidade de falha}
\simbololista{n_s}{\ensuremath{n_s}}{Número de amostras}
\simbololista{I_f}{\ensuremath{I_f}}{Função indicadora de falha}
\simbololista{g_lim}{\ensuremath{g}}{Função estado limite}
\simbololista{beta}{\ensuremath{\beta}}{Índice de confiabilidade}
\simbololista{Phi_inv}{\ensuremath{\Phi^{-1}}}{CDF inversa da normal padrão}
\simbololista{P_mat}{\ensuremath{P}}{Matriz de permutações}
\simbololista{R_mat}{\ensuremath{R}}{Matriz de números aleatórios}
\simbololista{S_mat}{\ensuremath{S}}{Matriz de amostras}
\simbololista{s_ij}{\ensuremath{s_{ij}}}{Elemento da matriz S}
\simbololista{F_inv_Xj}{\ensuremath{F^{-1}_{X_j}}}{Inversa da CDF da variável Xj}
\simbololista{x_ij}{\ensuremath{x_{ij}}}{Amostra da variável aleatória}
\simbololista{mu_g}{\ensuremath{\mu_g}}{Média da margem de segurança}
\simbololista{sigma_g}{\ensuremath{\sigma_g}}{Desvio padrão da margem de segurança}
\simbololista{mu_S}{\ensuremath{\mu_S}}{Média da solicitação}
\simbololista{sigma_S}{\ensuremath{\sigma_S}}{Desvio padrão da solicitação}
\simbololista{mu_R}{\ensuremath{\mu_R}}{Média da resistência}
\simbololista{sigma_R}{\ensuremath{\sigma_R}}{Desvio padrão da resistência}


%Para usar um dado símbolo SIMB ao longo do texto, use \gls{SIMB}.
